% Fragmento aditivo. No sustituye 10b_boltzmann.tex.
% Procedencia y distinción entre recuperación y formalización:
% 10c_boltzmann_procedencia.md. Etiquetas locales prefijadas ii-kb-aut.

\subsection{Mémoire généalogique, réduction informationnelle et transduction thermique}
\label{sec:ii-kb-aut-desarrollo}

La construction thermique utilise deux représentations du même développement
APP--TRIT--TPK. La première conserve les états, les chemins admissibles et
l’information du registre ; la seconde fournit l’action orientée
et la durée du retour. Leur composition détermine une énergie par
décision tritique. Les coordonnées de température et d’énergie sont introduites
ensuite au moyen d’applications dimensionnelles explicites. Cette séquence
permet d’identifier quelle information est conservée, quelle information extrait
un lecteur et quelle grandeur convertit cette information en entropie dimensionnelle.

\subsubsection{Lecteur, mémoire et récupération de l'état}

Fixons un niveau fini du développement et notons \(\Omega\) l'
ensemble des états enrichis admissibles. Le registre conserve feuillet,
orientation, résidu, quotient, retenue, route, phase et mémoire dans les
domaines où intervient chaque coordonnée. Soient
\[
 q:\Omega\longrightarrow Y,\qquad
 M:\Omega\longrightarrow\mathcal M
\]
l’application observable et la mémoire retenue, respectivement. Pour une
distribution \(p\) sur \(\Omega\), nous écrivons
\(\mathsf H(p)=-\sum_xp_x\log p_x\), avec \(0\log0=0\).
Le logarithme s’applique à des probabilités adimensionnelles.

\begin{definition}[Récupération généalogique à partir de l’observation et de sa mémoire]
Le couple \(\widehat q=(q,M)\) est un lecteur récupérable sur le support
\(\Omega_p\) lorsqu'il existe
\[
 d:\widehat q(\Omega_p)\longrightarrow\Omega_p,
 \qquad d\circ\widehat q=\operatorname{id}_{\Omega_p}.
\]
L'information de fibre de \(q\) est
\(\mathsf H_{\mathrm{fib}}(q;p)=\mathsf H(X\mid q(X))\), où
\(X\) a pour loi \(p\).
\end{definition}

\begin{proposition}[Identité de récupération et coût de la réduction]
Pour un lecteur récupérable,
\begin{align}
 \mathsf H(X\mid q(X),M(X))&=0,\\
 \mathsf H(X)&=\mathsf H(q(X))+\mathsf H(X\mid q(X)),\\
 \mathsf H(X\mid q(X))&=\mathsf H(M(X)\mid q(X)).
 \label{eq:ii-kb-aut-fibra}
\end{align}
\end{proposition}
\begin{proof}
L'application \(d\) restitue \(X\) à partir du couple \((q(X),M(X))\),
de sorte que la première entropie conditionnelle est nulle. Comme \(q\) est
une fonction de \(X\), la règle de la chaîne donne la seconde égalité.
Pour la troisième, on développe de deux manières
\(\mathsf H(X,M(X)\mid q(X))\). L'une fournit
\(\mathsf H(X\mid q(X))\), puisque \(M\) est une fonction de \(X\) ;
l'autre fournit
\(\mathsf H(M(X)\mid q(X))+\mathsf H(X\mid q(X),M(X))\).
\end{proof}

La division euclidienne des feuillets d'APP montre le mécanisme
arithmétique élémentaire. Avec \(i,j\in\{1,\ldots,9\}\), on conserve
\[
 i+j=9q^+_{ij}+R^+_{ij},\qquad
 ij=9q^\times_{ij}+R^\times_{ij},\qquad
 1\leq R^\sigma_{ij}\leq9\quad(\sigma\in\{+,\times\}).
\]
La paire résidu--quotient récupère la valeur entière de l'opération. En additionnant les différences sur les quatre-vingt-une positions, on obtient
\[
 2025-810=54+9\cdot129.
\]
La mémoire du quotient restaure la contribution que ne conserve pas la
réduction résiduelle. La récupérabilité de l’état complet nécessite,
en outre, de conserver les coordonnées de position et de chemin pertinentes ;
la paire résidu--quotient identifie ici le résultat de l’opération.

\subsubsection{Mise à jour réversible, décision tritique et effacement}

Soit \(U:\Omega\to\Omega\) la mise à jour TPK dans un domaine où le transport enrichi possède une inverse. La loi transportée est \(p'_y=p_{U^{-1}y}\). La bijection conserve le multiensemble des probabilités et, par conséquent,
\begin{equation}
 \mathsf H(U(X))=\mathsf H(X),\qquad
 \mathsf H(X\mid U(X))=0.
 \label{eq:ii-kb-aut-reversible}
\end{equation}
Si la sortie observée est \(q(U(X))\), l’information non observée
se quantifie par \(\mathsf H(X\mid q(U(X)))\). L’état complet,
sa lecture et la réinitialisation d’un registre sont ainsi trois applications distinctes.

Pour les trois états orientés du TRIT, soit
\(p=(p_{-1},p_0,p_{+1})\) une distribution normalisée. Son information est
\begin{equation}
 I_{\mathrm{TRIT}}(p)=-\sum_{a=-1}^{1}p_a\log p_a,
 \qquad 0\leq I_{\mathrm{TRIT}}(p)\leq\log3.
 \label{eq:ii-kb-aut-trit}
\end{equation}
L'extrémité supérieure correspond à \(p_a=1/3\). En effet,
\[
 D(p\Vert(1/3,1/3,1/3))
 =\sum_ap_a\log(3p_a)=\log3-I_{\mathrm{TRIT}}(p)\geq0;
\]
l'inégalité s'obtient à partir de la convexité de \(t\log t\), et l'égalité caractérise la distribution uniforme. L'extrémité inférieure correspond à une distribution concentrée sur un état.

La réalisation idéale de Landauer utilise une mémoire aux états logiques
énergétiquement dégénérés, un environnement isotherme à température \(\Theta\)
et un bilan qui conserve la sortie de l'application logique et comptabilise les
registres auxiliaires. Dans ce modèle, la diminution de l'entropie logique est
\(k_B[\mathsf H(X)-\mathsf H(F(X))]
=k_B\mathsf H(X\mid F(X))\) pour une application déterministe \(F\).
Le bilan entropique exige de transférer au milieu au moins cette quantité ;
son expression énergétique est \(Q\geq k_B\Theta\mathsf H(X\mid F(X))\).
La distinction entre évolution logique réversible et effacement physique
correspond à la réalisation de Landauer--Bennett \cite{bennett2003}.
La mise à jour inversible
\eqref{eq:ii-kb-aut-reversible} a une contribution logique nulle ; la
réinitialisation d'un triplet uniforme exige le transfert de \(\log3\) nats
au milieu et satisfait \(Q\geq k_B\Theta\log3\).
L'énergie de contrôle et la dissipation d'une réalisation physique
concrète appartiennent à son bilan thermodynamique supplémentaire.

\subsubsection{Action, fréquence et énergie de décision}

Soient \(\mathcal L_{\rm act}\) et \(\mathcal L_T\) les droites
dimensionnelles d'action et de durée ; la droite énergétique est
\(\mathcal L_E=\mathcal L_{\rm act}\otimes\mathcal L_T^*\).
Le développement précédent fournit les sections positives
\(\hbar_a\in\mathcal L_{\rm act}\), avec
\(a\in\{\mathrm{pre},\mathrm{ret}\}\), et la durée élémentaire
\(t_0\in\mathcal L_T\). Le calendrier complet fournit
\[
 T_{108}=108t_0,\qquad
 \omega_{108}=\frac{2\pi}{T_{108}},\qquad
 h_a=2\pi\hbar_a.
\]
La contraction dimensionnelle entre action et fréquence définit
\begin{equation}
 \epsilon_a:=\hbar_a\omega_{108}
 =\frac{h_a}{108t_0}
 =\frac{\hbar_a}{t_P},\qquad
 t_P=\frac{54}{\pi}t_0.
 \label{eq:ii-kb-aut-energia}
\end{equation}
Le dernier membre emploie la même période circulaire réduite que la section gravitationnelle. Les trois expressions correspondent à une unique section d'énergie, avec une orientation \(a\) maintenue dans tous les facteurs. L'information normalisée d'une décision uniforme produit la section d'énergie par nat
\begin{equation}
 \varepsilon_a:=\frac{\epsilon_a}{\log3}\in\mathcal L_E^+.
 \label{eq:ii-kb-aut-energia-nat}
\end{equation}
L'action précédemment construite détermine l'échelle d'énergie lorsque
la durée est fixée ; le cardinal des alternatives du TRIT détermine sa
normalisation informationnelle. La valeur de \(\log3\) intervient par le
dénombrement et la mesure déclarés dans \eqref{eq:ii-kb-aut-trit}.

\subsubsection{Définition dimensionnelle de Boltzmann et sélection des coordonnées}

Soit \(\mathcal L_\Theta\) la droite orientée de température. L'espace dimensionnel d'entropie est \(\mathcal L_{\rm ent}=\mathcal L_E\otimes\mathcal L_\Theta^*\). Une application linéaire positive \(k_B:\mathcal L_\Theta\to\mathcal L_E\) est, de manière équivalente, un élément positif de \(\mathcal L_{\rm ent}\). C'est pourquoi le même objet peut agir sur une température ou multiplier une information sans dimension.

\begin{definition}[Normalisation thermique de la décision tritique]
Un couple thermique compatible avec la section orientée \(a\) se compose d'
un transducteur positif \(k_B\) et d'une section positive
\(\Theta_{{\rm clk},a}\in\mathcal L_\Theta\) tels que
\begin{equation}
 k_B(\Theta_{{\rm clk},a})=\varepsilon_a,
 \qquad
 k_B\Theta_{{\rm clk},a}\log3=\epsilon_a.
 \label{eq:ii-kb-aut-par-termico}
\end{equation}
La constante de Boltzmann remplit la fonction de transducteur
énergie--température de l'information tritique dans cette réalisation.
\end{definition}

\Needspace{18\baselineskip}
\begin{samepage}
\begin{theorem}[Unicité relative à la section thermique et covariance dimensionnelle]
Une section positive \(\Theta_{{\rm clk},a}\) étant fixée, il existe une unique application linéaire positive qui satisfait \eqref{eq:ii-kb-aut-par-termico}. Pour des bases positives \(u_E,u_\Theta\), si
\[
 \epsilon_a=\epsilon_a^{\rm num}u_E,\quad
 \Theta_{{\rm clk},a}=\theta_a u_\Theta,\quad
 k_B(u_\Theta)=b\,u_E,
\]
sa coordonnée est
\begin{equation}
 b=\frac{\epsilon_a^{\rm num}}{\theta_a\log3}.
 \label{eq:ii-kb-aut-coordenada}
\end{equation}
Sous \(u'_E=s_Eu_E\), \(u'_\Theta=s_\Theta u_\Theta\), avec
\(s_E,s_\Theta>0\), on a
\[
 b'=\frac{s_\Theta}{s_E}b,\qquad
 \theta'_a=\frac{\theta_a}{s_\Theta},\qquad
 (\epsilon_a^{\rm num})'=\frac{\epsilon_a^{\rm num}}{s_E}.
\]
\end{theorem}
\end{samepage}
\begin{proof}
Une section non nulle engendre une droite. Prescrire son image détermine
l'unique application linéaire ; la positivité des deux sections prouve
sa positivité. Substituer les coordonnées dans
\eqref{eq:ii-kb-aut-par-termico} donne
\(b\theta_a\log3=\epsilon_a^{\rm num}\). Les trois lois de
transformation s'obtiennent en exprimant les mêmes sections et la même
application dans les nouvelles bases.
\end{proof}

Une même constante de Boltzmann doit agir dans les deux orientations.
Puisque \(t_P\) est commun, la compatibilité simultanée équivaut à
\begin{equation}
 \frac{\Theta_{{\rm clk},{\rm pre}}}{\Theta_{{\rm clk},{\rm ret}}}
 =\frac{\epsilon_{\rm pre}}{\epsilon_{\rm ret}}
 =\frac{\hbar_{\rm pre}}{\hbar_{\rm ret}}.
 \label{eq:ii-kb-aut-orientaciones}
\end{equation}
En effet, un isomorphisme de droites conserve les quotients de sections. Réciproquement, si l'égalité précédente est satisfaite, l'isomorphisme déterminé dans une orientation satisfait également la normalisation de l'autre. Ainsi, le retour modifie la section thermique du cycle dans la même proportion que l'énergie, tandis que le transducteur reste commun aux deux lectures.

L’équation précédente explicite le choix nécessaire pour fixer une
valeur individuelle. Si la section thermique n’a pas encore été sélectionnée
par un lecteur indépendant, pour chaque \(\lambda>0\), la paire
\[
 (k_B,\Theta_{{\rm clk},a})
 \longmapsto(\lambda k_B,\lambda^{-1}\Theta_{{\rm clk},a})
\]
conserve \(\epsilon_a\). Par conséquent, la présentation d’une
coordonnée individuelle sélectionnée en interne exige de donner le lecteur de
\(\Theta_{{\rm clk},a}\), les bases et leur transport dimensionnel.
Les équations d’action et d’information restent déterminées avant
cette spécification. L’identification ultérieure des bases au
joule et au kelvin appartient à l’expression métrologique du résultat.

\subsubsection{Entropie dimensionnelle et information de l'état réduit}

Pour une distribution de routes admises \(p\), ou pour un état réduit de matrice de densité \(\rho\), soient
\[
 s(p)=-\sum_xp_x\log p_x,\qquad
 s(\rho)=-\operatorname{Tr}(\rho\log\rho).
\]
Leurs expressions dimensionnelles sont
\begin{equation}
 S(p)=k_Bs(p),\qquad S(\rho)=k_Bs(\rho)
 \quad\text{en }\mathcal L_{\rm ent}.
 \label{eq:ii-kb-aut-entropia}
\end{equation}
L'évaluation en \(\Theta\) produit une énergie
\(\Theta S=k_B\Theta\,s\). Pour le triplet uniforme à la température
du cycle, \eqref{eq:ii-kb-aut-par-termico} donne exactement
\[
 S_{\rm TRIT}=k_B\log3,\qquad
 \Theta_{{\rm clk},a}S_{\rm TRIT}=\epsilon_a.
\]
Si \(\rho_A\) est la réduction de la purification bipartite construite dans la section surfacique, \(s(\rho_A)\) mesure l'intrication de cette bipartition. Le transport des unités conserve ses valeurs propres, son orientation et les occupations sectorielles. L'échelle surfacique \(\gamma a_0\ell_P^2\) convertit les occupations en aire ; le transducteur \(k_B\) convertit l'information en entropie dimensionnelle. Les deux lectures utilisent le même registre au moyen d'opérateurs distincts.

\subsubsection{Normalisation spectrale des histoires thermiques}

Soit \(\mathcal G\) un graphe fini d'extensions admissibles déjà
construit, de matrice d'adjacence \(A\). Chaque arête \(e\) conserve
ses coordonnées de route et reçoit un incrément adimensionnel d'action
\(a_*(e)>0\) du lecteur correspondant. Une fois choisie la section d'énergie
\(\epsilon_*\), posons \(t=\beta\epsilon_*\), où
\(\beta=(k_B\Theta)^{-1}\). L'opérateur thermique est
\begin{equation}
 (L_tf)(y)=\sum_{e:x\to y}\exp[-t a_*(e)]f(x).
 \label{eq:ii-kb-aut-operador}
\end{equation}
Tous ses exposants sont sans dimension. La condition
\(\rho(L_t)=1\) normalise la croissance pondérée des histoires.

\begin{proposition}[Existence et unicité de la pression normalisée]
Supposons que \(A\) est irréductible, \(\rho(A)>1\), et que
\(0<m\leq a_*(e)\leq M\) pour toutes les arêtes. Il existe un unique
\(t_*>0\) tel que \(\rho(L_{t_*})=1\), et
\begin{equation}
 \frac{\log\rho(A)}{M}\leq t_*\leq
 \frac{\log\rho(A)}{m}.
 \label{eq:ii-kb-aut-presion}
\end{equation}
\end{proposition}
\begin{proof}
La comparaison entrée par entrée fournit \(e^{-tM}A\leq L_t\leq e^{-tm}A\) pour \(t\geq0\). La monotonie du rayon spectral des matrices non négatives implique les bornes correspondantes pour \(\rho(L_t)\). Le rayon spectral est continu, vaut \(\rho(A)>1\) en zéro et tend vers zéro à l'infini. Pour \(s>0\), en outre, \(L_{t+s}\leq e^{-sm}L_t\), de sorte que \(\rho(L_{t+s})\leq e^{-sm}\rho(L_t)<\rho(L_t)\). La continuité et la décroissance stricte prouvent l'existence et l'unicité. Substituer la valeur un dans les bornes initiales démontre \eqref{eq:ii-kb-aut-presion}.
\end{proof}

La réalisation locale de trois alternatives équipondérées, représentée
par un sommet avec trois boucles et \(a_*(e)=1\), satisfait
\(\rho(L_t)=3e^{-t}\). Sa normalisation donne \(t_*=\log3\).
Avec \(\epsilon_*=\epsilon_a\), il en résulte
\(\beta_{\rm clk}=\log3/\epsilon_a\), exactement la relation
\eqref{eq:ii-kb-aut-par-termico}. Ce calcul identifie la
confluence locale tritique. Pour un graphe d'histoires avec des contraintes
supplémentaires, la confluence avec la même horloge exige de vérifier
\[
 \rho\!\left(L_{(\log3/\epsilon_a)\epsilon_*}\right)=1
\]
avec leurs incréments effectifs. La normalisation locale et la pression
du système complet conservent ainsi leurs domaines respectifs.

La relation obtenue réunit structure, valeur et sens. Le registre
détermine l’information récupérable ; les probabilités sur les
continuations déterminent sa lecture entropique ; l’action et la période
déterminent \(\epsilon_a\) ; le transducteur énergie--température
exprime l’entropie dimensionnelle. L’aire et l’intrication se
relient à travers l’état réduit et ses opérateurs de bord,
en maintenant visible la donnée qui produit chaque grandeur.
